% Source printed pp. 437-444 / PDF pp. 445-452.
\section{不稳定轨迹的 HKKN 分层}\label{src:8.7}
\begin{epigraph}
事实上, 就个人而言, 这一课题吸引我们的一个原因在于, 它将代数向量丛与 Morse 理论联系起来, 而我们各自在四分之一个世纪前对这两个课题作出过贡献\ldots 更妙的是, 促成这种相互作用的催化剂来自一个颇出人意料的领域, 即理论物理中的 Yang--Mills 方程.
\epigraphauthor{Michael Atiyah 与 Raoul Bott \cite{AB83}}
\end{epigraph}
对线性约化群 $G$ 在射影簇 $X\subseteq\mathbb P^n$ 上的作用, 我们将证明非半稳定轨迹 $X\setminus X^{\mathrm{ss}}$ 按照函数
\[
M(x):=\inf_{\lambda\in\mathbb X_*(G)}\frac{\mu(x,\lambda)}{\|\lambda\|}\in\mathbb R_{<0}
\]
分层为局部闭子概形, 称为\newterm{Hesselink--Kempf--Kirwan--Ness 分层}{Hesselink--Kempf--Kirwan--Ness stratification}, 简称 \textit{HKKN 分层}. 由 Kempf 最优不稳定化定理 (\bookref{8.6.6}), 对每个 $x\in X\setminus X^{\mathrm{ss}}$, 此下确界由某个最优不稳定化单参数子群 $\lambda$ 达到. 指数 $M(x)$ 越负, 点 $x$ 越非半稳定, 或越``不稳定''. 各层将由二元组 $(\lambda,M)$ 编号, 其中 $\lambda\in\mathbb X_*(G)$, $M\in\mathbb R_{<0}$. 如同 Białynicki--Birula 分层 (\bookref{7.8.14}), HKKN 分层是计算等变上同调的有力工具.

\subsection{Hesselink--Kempf--Kirwan--Ness 分层}\label{src:8.7.1-subsection}
对单参数子群 $\lambda:\mathbb G_m\to G$, 回顾 $X^\lambda:=\operatorname{Mor}^{\mathbb G_m}(\Spec\Bbbk,X)$ 是 $\lambda$-不动轨迹 (定义~\bookref{7.8.1}), $X_\lambda^+:=\operatorname{Mor}^{\mathbb G_m}(\mathbb A^1,X)$ 是参数化极限 $\lim_{t\to0}\lambda(t)\cdot x$ 存在的点 $x$ 的吸引轨迹 (定义~\bookref{7.8.6}). 用 $\operatorname{ev}_0:X_\lambda^+\to X^\lambda$ 表示由 $x\mapsto\lim_{t\to0}\lambda(t)\cdot x$ 定义的在 $0$ 处的取值映射. 若 $X$ 是射影的, Białynicki--Birula 分层 (\bookref{7.8.14}) 表明 $X_\lambda^+$ 是局部闭子概形的不交并. 还回顾 $C_\lambda$ 和 $P_\lambda$ 分别表示与 $\lambda$ 关联的中心化子和抛物子群 (定义~\bookref{B.1.27}). HKKN 分层对约化群也成立, 但与往常一样, 我们限制在线性约化群的情形, 因为我们仅在这一设置下发展了 GIT.
\begin{theorem}[HKKN 分层]\label{src:8.7.1}
设 $G$ 是代数闭域 $\Bbbk$ 上的线性约化群, 在 $\mathbb X_*(G)$ 上给定长度 $\|{-}\|$. 设 $X\subseteq\mathbb P(V)$ 是 $G$-等变闭子概形, $V$ 是有限维 $G$-表示. 对极大环面 $T\subseteq G$, 存在有限子集 $\Sigma\subseteq\mathbb X_*(T)\times\mathbb R_{<0}$, 以及将非半稳定轨迹分为 $G$-不变局部闭子概形的分层
\[
X\setminus X^{\mathrm{ss}}=\coprod_{(\lambda,M)\in\Sigma}S_{\lambda,M},
\]
使对每个 $(\lambda,M)\in\Sigma$ 有:
\begin{enumerate}[(1)]
\item 轨迹
\[
Y_{\lambda,M}:=\{x\in X_\lambda^+\mid M=M(x)=\mu(x,\lambda)/\|\lambda\|\}
\]
是 $X$ 的 $P_\lambda$-不变局部闭子概形, 且 $S_{\lambda,M}:=G\cdot Y_{\lambda,M}$ 是 $X$ 的 $G$-不变局部闭子概形;
\item 点 $x\in Y_{\lambda,M}$ 当且仅当 $\operatorname{ev}_0(x)=\lim_{t\to0}\lambda(t)\cdot x\in Y_{\lambda,M}\cap X^\lambda$; 因而
\[
Z_{\lambda,M}:=\{x\in X^\lambda\mid M(x)=M=\mu(x,\lambda)/\|\lambda\|\}
\]
是 $Y_{\lambda,M}$ 的 $C_\lambda$-不变闭子概形, 满足 $Y_{\lambda,M}=\operatorname{ev}_0^{-1}(Z_{\lambda,M})$;
\item 自然映射 $G\times^{P_\lambda}Y_{\lambda,M}\to S_{\lambda,M}$ 有限, 满射且普遍单射; 若 $\operatorname{char}(\Bbbk)=0$, 则此映射是同构;
\item 轨迹
\[
\bigcup_{\substack{(\lambda',M')\in\Sigma\\M'\leq M}}S_{\lambda',M'}
\]
闭, 特别地包含 $\overline S_{\lambda,M}$;
\item 若 $X$ 光滑, 则每个 $Y_{\lambda,M}$ 也光滑; 在 $\operatorname{char}(\Bbbk)=0$ 时, 各层 $S_{\lambda,M}$ 也光滑.
\end{enumerate}
\end{theorem}
\begin{remark}[层, 叶片与中心]\label{src:8.7.2}
轨迹 $S_{\lambda,M}$ 称为\newterm{层}{stratum}, $Y_\lambda$ 称为\newterm{叶片}{blade}, $Z_{\lambda,M}$ 称为\newterm{中心}{center} (有时称为\newterm{极限集}{limit set}). 用叠的语言, 存在分层
\[
[X/G]=[X^{\mathrm{ss}}/G]\amalg\coprod_{(\lambda,M)\in\Sigma}[S_{\lambda,M}/G],
\]
且对每个 $(\lambda,M)$ 存在图表
\begin{equation}\label{src:8.7.3}
\begin{tikzcd}[ampersand replacement=\&,column sep=large]
{[Z_{\lambda,M}/C_\lambda]}\arrow[r,"i"']\&{[Y_{\lambda,M}/P_\lambda]}\arrow[l,bend right=28,"\operatorname{ev}_0"']\arrow[r]\&{[X/G]}
\end{tikzcd}
\end{equation}
满足 $\operatorname{ev}_0\circ i=\mathrm{id}$. 在特征 $0$ 时, 存在等价 $[Y_{\lambda,M}/P_\lambda]\simeq[S_{\lambda,M}/G]$, 且映射 $[Y_{\lambda,M}/P_\lambda]\to[X/G]$ 是局部闭浸入.
\end{remark}
\begin{proof}[定理~\bookref{8.7.1} 的证明]
设 $\widehat X\subseteq\mathbb A(V)$ 是 $X$ 的仿射锥, $\widehat N\subseteq\widehat X$ 是零锥, 即 $X\setminus X^{\mathrm{ss}}$ 的仿射锥, 它以 $0\in\widehat N$ 为唯一闭 $G$-轨道. 对零锥 $\widehat N\subseteq\mathbb A(V)$ 应用命题~\bookref{8.6.16}, 得到有限子集 $\Sigma\subseteq\mathbb X_*(T)\times\mathbb R_{<0}$, 使对每个点 $\widehat x\in\widehat N\setminus0$, 存在唯一 $(\lambda,M)\in\Sigma$, 满足 $\lambda$ 是 $\widehat x$ 的最优不稳定化单参数子群且 $M=\mu(\widehat x,\lambda)/\|\lambda\|$. 因为 $\widehat N$ 仿射, 对每个 $(\lambda,M)\in\Sigma$, 轨迹 $\widehat N_\lambda^+\subseteq\widehat N$ 是闭子概形 (习题~\bookref{7.8.7}). 由于 $\widehat N_\lambda^+$ 在缩放下不变, 它对应于闭子概形 $Y_\lambda\subseteq X$. 考虑轨迹
\[
Y_{\lambda,M}\subseteq\widetilde Y_{\lambda,M}\subseteq Y_\lambda,
\]
其中 $\widetilde Y_\lambda=\{x\in Y_\lambda\mid M=\mu(x,\lambda)/\|\lambda\|\}$, $Y_{\lambda,M}:=\{x\in Y_\lambda\mid M=M(x)=\mu(x,\lambda)/\|\lambda\|\}$. 因为函数 $x\in Y_\lambda\mapsto\mu(x,\lambda)$ 局部常值, $\widetilde Y_\lambda\subseteq Y_\lambda$ 既开又闭. 使用等式
\[
Y_{\lambda,M}=\widetilde Y_{\lambda,M}\setminus\bigcup_{\substack{(\lambda',M')\in\Sigma\\M'<M}}\widetilde Y_{\lambda',M'},
\]
可知 $Y_{\lambda,M}\subseteq\widetilde Y_{\lambda,M}$ 开. 因为 $G$ 约化, $P_\lambda\subseteq G$ 是抛物子群, 且
\[
[Y_\lambda/P_\lambda]\simeq[G\times^{P_\lambda}Y_\lambda/G]\longrightarrow[X/G]
\]
是射影的, 所以 $\widetilde Y_{\lambda,M}$ 的像 $\widetilde S_{\lambda,M}=G\cdot\widetilde Y_{\lambda,m}\subseteq X$ 闭, $Y_{\lambda,M}$ 的像 $S_{\lambda,M}=G\cdot\widetilde S_{\lambda,M}\subseteq\widetilde S_{\lambda,M}$ 开. 这证明了 (1) 和 (4).

对于 (2), 若 $x\in X\subseteq\mathbb P(V)$, 则极限 $x_0=\lim_{t\to0}\lambda(t)\cdot x$ 是到子空间 $W=\bigoplus V_\chi$ 的投影, 其中直和遍历满足 $x$ 到 $V_\chi$ 的投影 $\operatorname{proj}_\chi(x)$ 非零且 $\langle\chi,\lambda\rangle=-\mu(x,\lambda)$ 的特征 $\chi\in\mathbb X^*(T)$. 由引理~\bookref{8.6.9}, $\lambda$ 位于闭凸集
\[
C_x=\{\lambda\in\mathbb X_*(T)_{\mathbb R}\mid\langle\chi,\lambda\rangle\geq1,\ \operatorname{proj}_\chi(x)\ne0\}
\]
中距原点最近的唯一点张成的射线上. 因而 $\lambda$ 也是类似定义的集合 $C_{x_0}$ 中距原点最近的点. 或者, 可以检验: 若 $\lambda_0\in\mathbb X_*(T)$ 是 $x_0$ 的最优不稳定化单参数子群, 则
\[
\frac{\mu(x_0,\lambda_0)}{\|\lambda_0\|}\leq\frac{\mu(x,\lambda)}{\|\lambda\|}
\]
(给出蕴涵 $x_0\in Y_{\lambda,M}\Rightarrow x\in Y_{\lambda,M}$), 且当 $N\gg0$ 时
\[
\frac{\mu(x,\lambda^N\lambda_0)}{\|\lambda^N\lambda_0\|}\leq\frac{\mu(x,\lambda)}{\|\lambda\|}
\]
(给出蕴涵 $x\in Y_{\lambda,M}\Rightarrow x_0\in Y_{\lambda,M}$).

对于 (3), 对 $x\in Y_{\lambda,M}$, 我们断言
\begin{equation}\label{src:8.7.4}
P_\lambda=\{g\in G(\Bbbk)\mid gx\in Y_{\lambda,M}\}.
\end{equation}
因为 $Y_{\lambda,M}$ 是 $P_\lambda$-不变的, 有包含关系 $\subseteq$. 反之, 若 $gx\in Y_{\lambda,M}$, 则 $\lambda$ 和 $g\lambda g^{-1}$ 都是 $x$ 的最优不稳定化单参数子群. 由 Kempf 最优不稳定化定理 (\bookref{8.6.5}), 抛物子群 $P_\lambda$ 与 $P_{g\lambda g^{-1}}=gP_\lambda g^{-1}$ 相等. 因为 $N_G(P_\lambda)=P_\lambda$ (命题~\bookref{B.1.29}), 得到 $g\in P_\lambda$. 由于 $[X_\lambda^+/P_\lambda]\to[X/G]$ 固有, $[Y_{\lambda,M}/P_\lambda]\to[S_{\lambda,M}/G]$ 也固有. 映射 $[Y_{\lambda,M}/P_\lambda]\to[S_{\lambda,M}/G]$ 按构造是满射, 由~\bookeqref{8.7.4} 在 $\Bbbk$-点上是单射; 因而它有限, 满射且普遍单射. 在 $\operatorname{char}(\Bbbk)=0$ 时, 它是同构.

对于 (5), 若 $X$ 光滑, 由定理~\bookref{7.8.9}, 每个 $Y_\lambda$ 也光滑. 因为 $Y_{\lambda,M}\subseteq Y_\lambda$ 开, $Y_{\lambda,M}$ 也光滑. 在 $\operatorname{char}(\Bbbk)=0$ 时, 由 (3), $S_{\lambda,M}=G\times^{P_\lambda}Y_{\lambda,M}$, 因而也光滑. 另见 \cite[\S3]{Hes81}, \cite[\S4]{Hes79} 及 \cite[定理 12.26 和 13.5]{Kir84}.
\end{proof}

\subsection{例子与相关分层}\label{src:8.7.2-subsection}
\begin{example}\label{src:8.7.5}
设 $\mathbb G_m$ 以权 $-1,2,3$ 线性作用在 $X=\mathbb P^2$ 上. 令 $\lambda=\mathrm{id}$ 为恒等单参数子群, 则非半稳定轨迹为 $V(x^2y,x^3z)$, 它具有分层 $S_{\lambda^{-1},-1}\cup S_{\lambda,-2}\cup S_{\lambda,-3}$, 其中
\begin{align*}
S_{\lambda^{-1},-1}&=\{[1:0:0]\},\\
S_{\lambda,-2}&=\{[0:y:z]\mid y\ne0\},\\
S_{\lambda,-3}&=\{[0:0:1]\}.
\end{align*}
\end{example}
\begin{example}\label{src:8.7.6}
重新考察带有 Segre 线性化的 $\mathrm{SL}_2$ 在 $X=(\mathbb P^1)^n$ 上的作用 (例~\bookref{8.6.11}). 设 $\lambda_0:\mathbb G_m\to\mathrm{SL}_2$ 由 $\lambda_0(t)=\operatorname{diag}(t^{-1},t)$ 定义. 若 $n$ 为奇数, 各层由 $(\lambda_0,-1),(\lambda_0,-3),\ldots,(\lambda_0,-n)$ 编号; 若 $n$ 为偶数, 则由 $(\lambda_0,-2),(\lambda_0,-4),\ldots,(\lambda_0,-n)$ 编号. 层 $S_{\lambda_0,n-2k}$ 由恰有 $k>n/2$ 个点重合的元组组成, 余维为 $k-1$. 叶片 $Y_{\lambda_0,n-2k}^+$ 由恰有 $k$ 个点为 $0$ 的元组组成, 而中心 $Z_{\lambda_0,n-2k}$ 是 $k$ 个点为 $0$ 且 $n-k$ 个点为 $\infty$ 的 $\mathbb G_m$-不动点集合.
\end{example}
\begin{remark}[Morse 分层]\label{src:8.7.7}
当 $X$ 是 $\mathbb C$ 上的光滑射影簇时, HKKN 分层与矩映射的平方范数 $\|{-}\|^2:X\to\mathbb R$ 的 Morse 分层相同. 给定 $x\in X$, 最优不稳定化单参数子群对应于从 $x$ 出发的最速下降路径. 中心 $Z_{\lambda,M}$ 对应于 $\|{-}\|^2$ 的临界值集合, 而各层 $S_{\lambda,M}$ 是由流向 $Z_{\lambda,M}$ 的点组成的局部闭子流形. 见 \cite[\S6]{Kir84} 和 \cite{Nes84}.
\end{remark}
\begin{remark}[\texorpdfstring{$\Theta$}{Theta}-分层]\label{src:8.7.8}
如注~\bookref{7.9.5} 所述, 存在等式
\[
\underline{\operatorname{Mor}}([\mathbb A^1/\mathbb G_m],[X/G])=\coprod_{\lambda\in\mathbb X_*(G)/\sim}[X_\lambda^+/P_\lambda],
\]
其中 $\mathbb X_*(G)/\sim$ 表示共轭意义下的单参数子群集合. 局部有限型代数 $\Bbbk$-叠 $\mathcal X$ 的一个\newterm{$\Theta$-分层}{$\Theta$-stratification} 由具有最小元素 $0\in\Sigma$ 的全序集 $\Sigma$ 及局部闭子叠的分层
\[
\mathcal X=\coprod_{\lambda\in\Sigma}\mathcal S_\lambda
\]
组成, 满足:
\begin{enumerate}[(1)]
\item 对每个 $\lambda\in\Sigma$, $\mathcal X_{\leq\lambda}:=\bigcup_{\rho\leq\lambda}\mathcal S_\rho$ 是 $\mathcal X$ 的开子叠;
\item 对每个 $\lambda\in\Sigma$, 存在连通分支的并 (称为 $\mathcal X_{\leq\lambda}$ 的一个 $\Theta$-层)
\[
\mathcal S_\lambda'\subseteq\underline{\operatorname{Mor}}([\mathbb A^1/\mathbb G_m],\mathcal X_{\leq\lambda}),
\]
使 $\operatorname{ev}_1:\mathcal S_\lambda'\to\mathcal X_{\leq\lambda}$ 是同构到 $\mathcal S_\lambda$ 上的闭浸入;
\item 对每个 $x\in|\mathcal X|$, 集合 $\{\lambda\in\Sigma\mid x\in|\mathcal X_{\leq\lambda}|\}$ 具有最小元素.
\end{enumerate}
见 \cite{HL14}. \textit{半稳定轨迹} $\mathcal X^{\mathrm{ss}}$ 按定义为开子叠 $\mathcal X_{\leq0}=\mathcal S_0$. 设 $\mathcal Z_\lambda'$ 是 $\mathcal S_\lambda'$ 在映射
\[
i:\underline{\operatorname{Mor}}(B\mathbb G_m,\mathcal X)\longrightarrow\underline{\operatorname{Mor}}([\mathbb A^1/\mathbb G_m],\mathcal X)
\]
下的逆像. 限制到 $0$ 所得的映射 $\operatorname{ev}_0:\underline{\operatorname{Mor}}([\mathbb A^1/\mathbb G_m],\mathcal X)\to\underline{\operatorname{Mor}}(B\mathbb G_m,\mathcal X)$ 是 $i$ 的截面, 且存在类似于~\bookeqref{8.7.3} 的图表
\[
\begin{tikzcd}[ampersand replacement=\&,column sep=large]
\mathcal Z_\lambda'\arrow[r,"i"']\&\mathcal S_\lambda'\arrow[l,bend right=28,"\operatorname{ev}_0"']\arrow[r,"\operatorname{ev}_1"]\&\mathcal X.
\end{tikzcd}
\]
在特征 $0$ 时, HKKN 分层是 $\Theta$-分层的一个例子, 其中先按 $-M$ 排列指标 $(\lambda,M)$, 再按 $\lambda$ 任意排序. 下一章将看到模叠 $\cBun_{r,d}(C)$ 具有称为 Harder--Narasimhan--Shatz 分层的 $\Theta$-分层 (见 \S\bookref{9.3.3-subsection}).
\end{remark}

\subsection{对上同调的应用: Kirwan 满射性}\label{src:8.7.3-subsection}
Kirwan 满射性 (\bookref{8.7.13}) 给出一个公式, 用 $X^{\mathrm{ss}}$ 及 HKKN 分层的中心 $Z_{\lambda,M}$ 的等变 Chow 群表示 $X$ 的等变 Chow 群. 为证明它, 需要一些预备引理.
\begin{lemma}\label{src:8.7.9}
设 $X$ 是代数闭域 $\Bbbk$ 上的光滑不可约概形, 带有光滑仿射代数群 $G$ 的作用. 设 $S_1,\ldots,S_r\subseteq X$ 是非空, 互不相交, 光滑, 不可约, 局部闭且 $G$-不变的子概形, 满足 $X=\coprod_iS_i$, 且对每个 $i$, $S_{\geq i}:=\bigcup_{j\geq i}S_j$ 闭. 设 $d_i$ 是 $S_i$ 在 $X$ 中的余维. 若对每个 $i$, 最高 Chern 类 $c_{d_i}^G(N_{S_i/X})\in\mathrm{CH}_G^*(S_i)_{\mathbb Q}$ 是非零因子, 则对每个 $k$ 有
\[
\dim\mathrm{CH}_G^k(X)_{\mathbb Q}=\sum_{i=1}^r\dim\mathrm{CH}_G^{k-d_i}(S_i)_{\mathbb Q}.
\]
\end{lemma}
\begin{proof}
由假设, 对每个 $i$, $S_{\leq i}=\bigcup_{j\leq i}S_j$ 开, 且 $S_i\subseteq S_{\leq i}$ 是具有开补集 $S_{<i}$ 的闭子概形. 有交换图
\[
\begin{tikzcd}[ampersand replacement=\&,column sep=small]
\mathrm{CH}_G^{k-d_i}(S_i)\arrow[r]\arrow[dr]\&\mathrm{CH}_G^k(S_{\leq i})\arrow[r]\arrow[d]\&\mathrm{CH}_G^k(S_{<i})\arrow[r]\&0\\
\&\mathrm{CH}_G^k(S_i)\&\&
\end{tikzcd}
\]
其中上行是右正合的切除序列 (\bookref{7.1.36}(3)), 竖直向下的箭头由与 $S_i$ 求交给出. 由自交公式 (\bookref{7.1.36}(6)), 复合 $\mathrm{CH}_G^{k-d_i}(S_i)\to\mathrm{CH}_G^k(S_i)$ 是乘以 $c_{d_i}^G(N_{S_i/X})$. 由假设, 与 $\mathbb Q$ 作张量积后此映射为单射. 因而上行与 $\mathbb Q$ 作张量积后成为正合序列, 且
\[
\dim\mathrm{CH}_G^k(S_{\leq i})_{\mathbb Q}=\dim\mathrm{CH}_G^{k-d_i}(S_i)_{\mathbb Q}+\dim\mathrm{CH}_G^k(S_{<i})_{\mathbb Q}.
\]
所求公式由归纳得到. 另见 \cite[命题 1.9]{AB83}.
\end{proof}
\begin{remark}\label{src:8.7.10}
若 $[S_i/G]$ 是 Deligne--Mumford 叠, 则当 $k\gg0$ 时 $\mathrm{CH}_G^k(S_i)$ 消失, 且 $c_{d_i}^G(N_{S_i/X})$ 是零因子.
\end{remark}
下面给出最高 Chern 类为非零因子的一个便利判据.
\begin{lemma}\label{src:8.7.11}
设 $X$ 是代数闭域 $\Bbbk$ 上的光滑不可约概形, 带有连通, 光滑, 仿射代数群 $G$ 的作用, $N$ 是 $X$ 上秩为 $d$ 的 $G$-等变向量丛. 假定存在子群 $\mathbb G_m\subseteq G$ 平凡作用在 $X$ 上, 且存在点 $x\in X(\Bbbk)$ 使 $N\otimes\kappa(x)$ 不含 $\mathbb G_m$-不变向量. 则 $c_d^G(N)\in\mathrm{CH}_G^*(X)_{\mathbb Q}$ 是非零因子.
\end{lemma}
\begin{proof}
取包含 $\mathbb G_m$ 的极大环面 $T$, 以及特征 $T\to\mathbb G_m$, 使复合 $\mathbb G_m\hookrightarrow T\to\mathbb G_m$ 由 $t\mapsto t^d$, $d>0$, 给出. 由~\bookref{7.1.38}(b), $\mathrm{CH}_G^*(X)_{\mathbb Q}=\mathrm{CH}_T^*(X)_{\mathbb Q}^W$, 其中 $W$ 是 Weyl 群. 因为 $\mathrm{CH}_G^*(X)_{\mathbb Q}$ 是 $\mathrm{CH}_T^*(X)_{\mathbb Q}$ 的子环, 归约为证明 $c_d^T(N)\in\mathrm{CH}_T^*(X)_{\mathbb Q}$ 是非零因子. 若将 $T$ 写成给定的 $\mathbb G_m$ 与子环面 $T'$ 的乘积, 则由~\bookref{7.1.38}(a),
\[
\mathrm{CH}_T^*(X)\simeq\mathrm{CH}_{T'}^*(X)\otimes\mathrm{CH}^*(B\mathbb G_m)\simeq\mathrm{CH}_{T'}^*(X)[z].
\]
对 $x\in X(\Bbbk)$, 可写成
\begin{align*}
c_d^T(N)&=\sum c_i^{T'}(N)\otimes c_{d-i}^{\mathbb G_m}(N\otimes\kappa(x))\\
&=1\otimes c_d^{\mathbb G_m}(N\otimes\kappa(x))+\text{更高次项}.
\end{align*}
若 $a_1,\ldots,a_d$ 表示 $N\otimes\kappa(x)$ 的 $\mathbb G_m$-权, 则由假设, 每个 $a_i\ne0$, 且
\[
c_d^{\mathbb G_m}(N\otimes\kappa(x))=\Big(\prod_i a_i\Big)z^d\in\mathrm{CH}^*(B\mathbb G_m)_{\mathbb Q}\simeq\mathbb Q[z]
\]
是非零因子, 因而 $c_d^T(N)$ 也是非零因子. 另见 \cite[命题 13.4]{AB83} 及 \cite[\S3.2]{Bri97}.
\end{proof}
定义 $G$-等变概形 $X$ 的\newterm{$G$-等变 Chow--Poincaré 多项式}{$G$-equivariant Chow--Poincaré polynomial} 为
\[
p_G(X,t)=\sum_{d=0}^\infty\big(\dim\mathrm{CH}_G^d(X)_{\mathbb Q}\big)t^d.
\]
也用 $p(X,t)=\sum_{d=0}^\infty\big(\dim\mathrm{CH}^d(X)_{\mathbb Q}\big)t^d$ 表示 (非等变的) \textit{Chow--Poincaré 多项式}.
\begin{exercise}\label{src:8.7.12}
设 $X$ 是代数闭域 $\Bbbk$ 上的光滑, 不可约, 拟射影概形, 带有 $\mathbb G_m$ 的作用, 且 $X^+\to X$ 满射 (即 $X$ 是射影的). 设 $X=\coprod_{i=1}^rX_i$ 和 $X^{\mathbb G_m}=\coprod_{i=1}^rF_i$ 是 Białynicki--Birula 分层 (\bookref{7.8.14}), $d_i$ 是 $X_i$ 在 $X$ 中的余维. 证明
\[
p_{\mathbb G_m}(X,t)=\sum_{i=1}^rp(F_i,t)\cdot t^{d_i}(1-t)^{-1}.
\]
\end{exercise}
\begin{proposition}[Kirwan 满射性]\label{src:8.7.13}
在 HKKN 分层 (\bookref{8.7.1}) 的假设下, 进一步假定 $X$ 光滑且不可约, $\operatorname{char}(\Bbbk)=0$. 假定对所有 $(\lambda,M)$, 层 $S_{\lambda,M}$ 是余维为 $d_{\lambda,M}$ 的等维概形. 则
\[
\dim\mathrm{CH}_G^k(X)_{\mathbb Q}=\dim\mathrm{CH}_G^k(X^{\mathrm{ss}})_{\mathbb Q}+\sum_{(\lambda,M)}\dim\mathrm{CH}_{C_\lambda}^{k-d_{\lambda,M}}(Z_{\lambda,M})_{\mathbb Q},
\]
且
\[
p_G(X,t)=p_G(X^{\mathrm{ss}},t)+\sum_{(\lambda,M)}p_{C_\lambda}(Z_{\lambda,M},t)t^{d_{\lambda,M}}.
\]
\end{proposition}
\begin{proof}
由 HKKN 分层 (\bookref{8.7.1}), $[S_{\lambda,M}/G]\simeq[Y_{\lambda,M}/P_\lambda]$; 由 Białynicki--Birula 分层 (\bookref{7.8.14}), $\operatorname{ev}_0:Y_{\lambda,M}\to Z_{\lambda,M}$ 是 Zariski 局部仿射纤维化, 且关于 $P_\lambda\to C_\lambda$ 等变. 我们断言 $[Y_{\lambda,M}/P_\lambda]\to[Z_{\lambda,M}/C_\lambda]$ 诱导同构
\begin{equation}\label{src:8.7.14}
\mathrm{CH}_{C_\lambda}^*(Z_{\lambda,M})\longrightarrow\mathrm{CH}_{P_\lambda}^*(Y_{\lambda,M}).
\end{equation}
由等变 Chow 群的定义 (\S\bookref{7.1.6-subsection}), $\mathrm{CH}_{P_\lambda}^i(Y_{\lambda,M})$ 等同于 $\mathrm{CH}^i(Y_{\lambda,M}\times^{P_\lambda}V)$, 其中 $V$ 是某个 $P_\lambda$-表示 $W$ 的 $\mathbb A(W)$ 中的开子空间, 满足 $P_\lambda$ 在 $V$ 上自由作用, 且 $\mathbb A(W)\setminus V$ 的余维充分大. 类似地, $\mathrm{CH}_{C_\lambda}^i(Z_{\lambda,M})$ 等同于 $\mathrm{CH}^i(Z_{\lambda,M}\times^{C_\lambda}V)$. 映射~\bookeqref{8.7.14} 对应于由复合
\[
Y_{\lambda,M}\times^{P_\lambda}V\longrightarrow Z_{\lambda,M}\times^{P_\lambda}V\longrightarrow Z_{\lambda,M}\times^{C_\lambda}V
\]
诱导的 Chow 群上的拉回映射. 第一个映射是 Zariski 局部仿射纤维化, 第二个映射是 $U_\lambda=\ker(P_\lambda\to C_\lambda)$ 下的主丛. 因为 $U_\lambda$ 幺幂, 它同构于仿射空间 (见例~\bookref{B.1.23}), 且主 $U_\lambda$-丛在 Zariski 拓扑中局部平凡 (见例~\bookref{B.1.61}). 由性质~\bookref{7.1.36}(2), 得到~\bookeqref{8.7.14} 是同构.

我们还断言 $c_{d_{\lambda,M}}(N_{S_{\lambda,M}/X})\in\mathrm{CH}_G^*(S_{\lambda,M})$ 是非零因子. 因为在 $\mathrm{CH}_C^*(Z_{\lambda,M})\simeq\mathrm{CH}_G^*(S_{\lambda,M})$ 下, $(N_{S_{\lambda,M}/X})|_{Z_{\lambda,M}}$ 与 $N_{S_{\lambda,M}/X}$ 等同, 只需证明
\[
c_{d_{\lambda,M}}((N_{S_{\lambda,M}/X})|_{Z_{\lambda,M}})\in\mathrm{CH}_C^*(Z_{\lambda,M})_{\mathbb Q}
\]
是非零因子. 由 Białynicki--Birula 分层 (\bookref{7.8.14}), $\lambda$ 在法丛的一个纤维上以非零权作用. 因而引理~\bookref{8.7.11} 表明 $c_{d_{\lambda,M}}((N_{S_{\lambda,M}/X})|_{Z_{\lambda,M}})$ 是非零因子.

于是可应用引理~\bookref{8.7.9}, 将各层 $S_{\lambda,M}$ 先按 $-M$ 排序, 再对 $\lambda$ 任意排序; 半稳定轨迹 $U^{\mathrm{ss}}$ 视为具有最小指标的层. 得到
\begin{align*}
\dim\mathrm{CH}_G^k(X)_{\mathbb Q}
&=\dim\mathrm{CH}_G^k(X^{\mathrm{ss}})_{\mathbb Q}+\sum_{(\lambda,M)}\dim\mathrm{CH}_G^{k-d_{\lambda,M}}(S_{\lambda,M})_{\mathbb Q}\\
&=\dim\mathrm{CH}_G^k(X^{\mathrm{ss}})_{\mathbb Q}+\sum_{(\lambda,M)}\dim\mathrm{CH}_{C_{\lambda,M}}^{k-d_{\lambda,M}}(Z_{\lambda,M})_{\mathbb Q},
\end{align*}
Chow--Poincaré 多项式的公式也随之得到.
\end{proof}
\begin{remark}\label{src:8.7.15}
\cite[定理 5.4]{Kir84} 对 de Rham 上同调建立了此公式. 此时使用 Thom--Gysin 长正合序列
\[
\cdots\to\mathrm H_G^{k-d_{\lambda,M}}(S_{\lambda,M})\to\mathrm H_G^k(S_{\leq(\lambda,M)})\to\mathrm H_G^k(S_{<(\lambda,M)})\to\cdots
\]
代替切除序列
\[
\mathrm{CH}_G^{k-d_{\lambda,M}}(S_{\lambda,M})\to\mathrm{CH}_G^k(S_{\leq(\lambda,M)})\to\mathrm{CH}_G^k(S_{<(\lambda,M)})\to0.
\]
Kirwan 通过证明右侧每个映射均满射, 证明长正合序列分裂为短正合序列, 这解释了为何称此结果为 \textit{Kirwan 满射性}. 为此, 注意对所有 $(\lambda,M)$ 右侧映射满射, 等价于对所有 $(\lambda,M)$ 左侧映射单射; 后一条件可与上面一样, 通过论证法丛的最高 Chern 类是非零因子来验证.
\end{remark}
\begin{example}\label{src:8.7.16}
使用例~\bookref{8.7.6} 确定的 HKKN 分层, 计算 $[(\mathbb P^1)^{n,\mathrm{ss}}/\mathrm{SL}_2]$ 的有理 Chow 群的维数. 当 $n$ 为奇数时, 由性质~\bookref{7.1.36}(4), 这也给出 GIT 商 $(\mathbb P^1)^{n,\mathrm{ss}}/\mathrm{SL}_2$ 的有理 Chow 群的维数.

因为 $[(\mathbb P^1)^n/\mathrm{SL}_2]\to B\mathrm{SL}_2$ 是迭代的 $\mathbb P^1$-丛, 且 $\mathrm{CH}^*(\mathrm{SL}_2)\simeq\mathbb Z[T]$ 由次数 $2$ 的元素生成,
\begin{align*}
\mathrm{CH}^*([(\mathbb P^1)^n/\mathrm{SL}_2])
&\simeq\mathrm{CH}^*((\mathbb P^1)^n)\otimes\mathrm{CH}^*(B\mathrm{SL}_2)\\
&\simeq\mathbb Z[H_1,\ldots,H_n]/(H_1,\ldots,H_n)^2\otimes\mathbb Z[T],
\end{align*}
Chow--Poincaré 多项式为 $p_{\mathrm{SL}_2}((\mathbb P^1)^n,t)=(1+t)^n(1-t^2)^{-1}$. 另一方面, 恰有 $k$ 个点相同的层 $S_{\lambda,n-2k}$ 余维为 $k-1$, 其中心 $Z_{\lambda,n-2k}$ 由 $\binom nk$ 个 $\mathbb G_m$-不动点组成. 因而 $p_{\mathbb G_m}(Z_{\lambda,n-2k},t)=\binom nk(1-t)^{-1}$, 且
\begin{align*}
p_{\mathrm{SL}_2}((\mathbb P^1)^{\mathrm{ss}},t)
&=(1+t)^n(1-t)^{-1}-\sum_{k>n/2}\binom nk t^{k-1}(1-t)^{-1}\\
&=1+nt+\cdots\\
&\quad+\left(1+(n-1)+\binom{n-1}2+\cdots+\binom{n-1}{\min(d,n-3-d)}\right)t^d\\
&\quad+\cdots+nt^{n-4}+t^{n-3}.
\end{align*}
另见 \cite[\S16.1]{Kir84}.
\end{example}
